\section{Additional Results}
\label{sec:app-additional}

\FloatBarrier
\subsection{Key Invariance (E03)}

To verify that distinguishers learn cipher-structural features rather than key-specific patterns, we evaluate a trained \speck 5r model on data encrypted with permuted keys (keys with shuffled bytes).
The model scores 50.01\% ($\pm$0.16)---indistinguishable from random---confirming that the learned features are key-invariant.
When evaluated on the original key distribution, the same model achieves 87.10\%.

\FloatBarrier
\subsection{Robustness to Noise (E04)}

We corrupt test ciphertext bits with Gaussian noise or random bitflips and measure accuracy degradation on \speck 5r.

\begin{table}[htbp]
\centering
\caption{Accuracy under input perturbation (\speck 5r, baseline: 87.10\%).}
\label{tab:robustness}
\begin{tabular}{@{}lcclcc@{}}
\toprule
\multicolumn{3}{c}{\textbf{Gaussian Noise ($\sigma$)}} & \multicolumn{3}{c}{\textbf{Bitflip Rate}} \\
\cmidrule(r){1-3}\cmidrule(l){4-6}
$\sigma$ & Acc.\ (\%) & Drop & Rate & Acc.\ (\%) & Drop \\
\midrule
0.01 & 87.10 & 0.0  & 0.01 & 83.97 & $-3.1$ \\
0.05 & 86.90 & $-0.2$ & 0.05 & 73.87 & $-13.2$ \\
0.10 & 85.62 & $-1.5$ & 0.10 & 65.33 & $-21.8$ \\
0.20 & 75.05 & $-12.1$ & 0.15 & 59.68 & $-27.4$ \\
0.50 & 53.41 & $-33.7$ & 0.20 & 56.17 & $-31.0$ \\
1.00 & 50.77 & $-36.3$ & 0.30 & 52.09 & $-35.0$ \\
\bottomrule
\end{tabular}
\end{table}

The network degrades gracefully under Gaussian noise ($\sigma \leq 0.10$: $<$2\pp loss) but is highly sensitive to bitflips (1\% flip rate: $-3.1$\pp).
This asymmetry reflects the discrete, bit-level nature of the learned features.

\FloatBarrier
\subsection{Markov Depth Analysis (E05)}

We train distinguishers using different numbers of intermediate round outputs as input features (``Markov depth'') to measure how much historical state information helps.

\begin{table}[htbp]
\centering
\caption{Accuracy vs.\ Markov depth on \speck (predicting round 7 label).}
\label{tab:markov-depth}
\begin{tabular}{@{}rc@{}}
\toprule
\textbf{Depth (rounds visible)} & \textbf{Accuracy (\%)} \\
\midrule
1 & $52.07 \pm 0.52$ \\
2 & $67.67 \pm 0.91$ \\
3 & $87.32 \pm 0.36$ \\
5 & $99.95 \pm 0.03$ \\
7 & $100.00 \pm 0.00$ \\
\bottomrule
\end{tabular}
\end{table}

Performance saturates at depth 3 (87.3\%), consistent with the last 3 rounds containing the bulk of exploitable differential information.

\FloatBarrier
\subsection{Conditional Mutual Information (E06)}

We estimate mutual information between round labels and ciphertext difference, conditioned on the previous round's output.
Conditional MI is near-zero across all rounds ($\leq 4.5 \times 10^{-5}$ bits), confirming that round transitions approximately satisfy the Markov property---conditioning on the current round state captures all information relevant to the next round.

\begin{figure}[htbp]
  \centering
  \includegraphics[width=\columnwidth]{fig7_signal_decay.png}
  \caption{Signal decay heatmap for \speck: accuracy across 7 input differences and 7 round counts (2--8).
  All differences produce similar accuracy at a given round, with the default 0x00400000 showing the strongest signal overall.
  The green-to-red gradient visually tracks the transition from distinguishable to random.}
  \label{fig:signal-decay}
\end{figure}

The heatmap in \Cref{fig:signal-decay} shows a $7 \times 7$ grid of distinguishing accuracy (7 input differences $\times$ 7 round counts, 2--8).
All differences produce similar accuracy at the same round count, with variation $<$3\pp.
The default difference (0x00400000) is not optimal at low rounds: 0x00100000 achieves comparable accuracy with more consistent cross-round behavior.

\FloatBarrier
\subsection{Input Difference Search (E10)}

\begin{table}[htbp]
\centering
\caption{Top-5 single-bit input differences for \speck 5r.}
\label{tab:diff-search}
\begin{tabular}{@{}lcc@{}}
\toprule
\textbf{$\Delta P$} & \textbf{Accuracy (\%)} & \textbf{Bit position} \\
\midrule
\texttt{0x00008000} & $76.20 \pm 0.80$ & 15 (MSB low word) \\
\texttt{0x00004000} & $70.87 \pm 1.04$ & 14 \\
\texttt{0x00020000} & $68.73 \pm 1.42$ & 17 \\
\texttt{0x00040000} & $68.28 \pm 1.70$ & 18 \\
\texttt{0x00010000} & $67.69 \pm 1.75$ & 16 \\
\bottomrule
\end{tabular}
\end{table}

The optimal single-bit difference (0x00008000, bit 15) achieves 76.2\% at 5 rounds, 12\pp below Gohr's known-good difference (0x00400000, 88.7\% with R2 representation) but using no prior cryptanalytic knowledge.
The ranking shows clear positional dependence: bits 14--18 (near the modular addition boundary) dominate, connecting to the carry propagation structure of SPECK's ARX operations.
This complements Bellini et al.'s~\citep{bellini2023dbitnet} evolutionary search with a simpler exhaustive analysis over single-bit differences.

\FloatBarrier
\subsection{Generative Markov Validation (E16)}

\begin{table}[htbp]
\centering
\caption{Markov vs.\ memory VAE prediction on \speck round transitions.
         Ratio $\approx 1.0$ indicates Markov property holds.}
\label{tab:markov-vae}
\begin{tabular}{@{}rccc@{}}
\toprule
\textbf{Round} & \textbf{Markov MSE} & \textbf{Memory MSE} & \textbf{Ratio} \\
\midrule
1 & 0.1060 & 0.1060 & $1.000 \pm 0.000$ \\
2 & 0.1593 & 0.1592 & $1.000 \pm 0.000$ \\
3 & 0.1879 & 0.1873 & $0.997 \pm 0.000$ \\
4 & 0.1912 & 0.1909 & $0.999 \pm 0.001$ \\
5 & 0.2007 & 0.2027 & $1.010 \pm 0.003$ \\
6 & 0.2000 & 0.2075 & $1.037 \pm 0.018$ \\
\bottomrule
\end{tabular}
\end{table}

\begin{figure}[htbp]
  \centering
  \includegraphics[width=\columnwidth]{fig8_vae_prediction.png}
  \caption{Generative Markov validation (E16).
  \emph{Left:} Reconstruction MSE for Markov (single-round input) vs.\ Memory (two-round input) VAEs across round transitions.
  \emph{Right:} Memory advantage ratio---values $>$1.0 indicate the memory model performs \emph{worse}, confirming the Markov property holds through round 4 and weakly breaks at rounds 5--6.}
  \label{fig:vae-prediction}
\end{figure}

\begin{figure*}[htbp]
  \centering
  \includegraphics[width=\textwidth]{fig9_vae_latent.png}
  \caption{VAE latent space visualizations for \speck across rounds 1--6.
  Red: cipher-generated ciphertext pairs; grey: random pairs.
  At round 1, cipher and random distributions are cleanly separated.
  By round 4, substantial overlap emerges.
  At round 6, the distributions are nearly indistinguishable, visually confirming the diminishing differential signal at higher rounds.}
  \label{fig:vae-latent}
\end{figure*}

The Markov property holds tightly through round 4 (ratio within 0.3\% of 1.0) but drifts at rounds 5--6 (ratio up to 1.037; \Cref{tab:markov-vae,fig:vae-prediction}), suggesting that deep-round transitions exhibit weak memory effects.
We note this honestly: the claim ``Markov ratio $1.000 \pm 0.004$'' in the original manuscript was misleading---it applied only to rounds 1--4.

The VAE latent space (\Cref{fig:vae-latent}) provides a complementary visualization: at early rounds, the cipher and random distributions occupy distinct regions of the learned manifold; by round 6, they merge completely, consistent with the distinguisher's accuracy collapsing to 50\%.

\FloatBarrier
\subsection{RL-Based Differential Search (E18)}

We train a reinforcement learning agent to discover good input differences for \speck.
At 5 rounds, RL fails: the best RL-discovered difference achieves 50.6\% (vs.\ 55.0\% for the known difference), indistinguishable from the random baseline.
At 3 rounds, RL shows partial success: one seed discovers 0x00102000 achieving 85.0\% (vs.\ 86.6\% for the known difference), but the other seeds find only random-quality differences.
The RL approach's reward signal becomes too noisy at higher rounds for the agent to converge, suggesting that gradient-free search for input differences requires substantially more samples than gradient-based training of the distinguisher itself.

\FloatBarrier
\subsection{SIMON Cross-Round Transfer (Full Table)}

\begin{table}[htbp]
\centering
\caption{Full cross-round transfer for \simon (source: 6r, accuracy 98.1\%).}
\label{tab:simon-transfer-full}
\begin{tabular}{@{}lccc@{}}
\toprule
\textbf{Target} & \textbf{Accuracy (\%)} & \textbf{$p$-value} & \textbf{Direction} \\
\midrule
$\to$ 4r & $48.61 \pm 0.10$ & $9.4 \times 10^{-6}$ & Anti-transfer \\
$\to$ 5r & $48.83 \pm 0.34$ & $0.002$ & Anti-transfer \\
$\to$ 7r & $50.27 \pm 0.11$ & $0.009$ & Weak positive \\
$\to$ 8r & $50.96 \pm 0.10$ & $4.3 \times 10^{-5}$ & Weak positive \\
$\to$ 9r & $50.09 \pm 0.12$ & $0.211$ & n.s. \\
\bottomrule
\end{tabular}
\end{table}

\simon exhibits anti-transfer for $r' < r$ (same pattern as \speck) and weak positive transfer for $r' = r+1$.
The asymmetry is informative: features learned at 6 rounds partially apply to 7--8 rounds (positive direction) but actively mislead at 4--5 rounds (negative direction).
